\section{Relating the Sigmoid Inflection Point to Mode Count}\label{app:k0-derivation}

\textbf{Note:} The derivation below assumes that modes are approximately independent (seeing a response from mode $i$ does not help $\theta$ predict mode $j$).
Section~\ref{sec:cross-mode-learning} shows this assumption fails for capable base models, which exhibit cross-mode learning and produce exponential decay at all $m/n$ ratios rather than the predicted sigmoid.
The derivation is included for theoretical completeness and because it does hold for weaker models (GPT-2).

\subsection{Setup}

We have $m$ modes with uniform weights $w_i = 1/m$.
Each of $n$ responses is drawn independently from one of the $m$ modes.
We fit a sigmoid to the permutation-averaged progressive conditional surprise curve:
\begin{equation}
    \bar{a}_k = a_\infty + \frac{\alpha}{1 + e^{\beta(k - k_0)}}
\end{equation}
where $k_0$ is the inflection point, $\alpha = \bar{a}_1 - a_\infty$ is the total learnable gap, $\beta > 0$ controls the steepness of the transition, and $a_\infty$ is the asymptotic floor.
We seek the relationship between $k_0$ and $m$.

\subsection{The Mode Discovery Crossover}

After $k$ responses drawn uniformly from $m$ modes, the expected probability that the next response comes from a mode not yet observed is:
\begin{equation}
    p_{\mathrm{new}}(k) = \left(1 - \frac{1}{m}\right)^{k}
\end{equation}

This decreases monotonically from $p_{\mathrm{new}}(0) = 1$ toward 0.
The complementary probability $1 - p_{\mathrm{new}}(k)$ is the chance that the next response comes from an already-seen mode.

Define $k^*$ as the \textbf{discovery crossover}: the number of responses at which a new response is equally likely to come from a novel mode or a previously seen mode.
Setting $p_{\mathrm{new}}(k^*) = 1/2$:
\begin{equation}
    \left(1 - \frac{1}{m}\right)^{k^*} = \frac{1}{2}
\end{equation}
\begin{equation}\label{eq:kstar}
    k^* = \frac{\ln 2}{-\ln(1 - 1/m)} = \frac{\ln 2}{\ln\!\left(\frac{m}{m-1}\right)}
\end{equation}

For large $m$, using $\ln(1 - 1/m) \approx -1/m$:
\begin{equation}\label{eq:kstar-approx}
    k^* \approx m \ln 2 \approx 0.693\, m
\end{equation}

This is an exact combinatorial quantity.
Before $k^*$, mode discovery dominates: most new responses reveal modes $\theta$ has not yet seen.
After $k^*$, mode refinement dominates: most new responses come from modes $\theta$ has already encountered.

For reference, the expected number of distinct modes observed at $k^*$ is:
\begin{equation}
    \E[\text{distinct modes at } k^*] = m\!\left(1 - \left(1 - \frac{1}{m}\right)^{k^*}\right) = m \cdot \frac{1}{2} = \frac{m}{2}
\end{equation}

That is, at the crossover, exactly half the modes have been seen in expectation.

\subsection{Hypothesis: \texorpdfstring{$k_0 \approx k^*$}{k0 approx k*}}

We hypothesize that the sigmoid inflection point $k_0$ is near the discovery crossover $k^*$.
The reasoning is as follows.
The inflection point of the learning curve is where $\theta$'s marginal learning rate is greatest.
Before $k^*$, $\theta$ is still in the discovery-dominated regime: it encounters new modes frequently, but has seen too few examples of any single mode to fully learn within-mode structure.
After $k^*$, $\theta$ is in the refinement-dominated regime: it rarely encounters new modes, and each additional example of a known mode provides diminishing returns.
The learning rate should peak near the transition, where $\theta$ has seen enough modes to start forming generalizations and enough repeats to start exploiting them.

This is a qualitative argument, not a proof.
The true inflection point depends on how much information $\theta$ extracts from a mode discovery versus a mode refinement, which is a property of $\theta$ and the mode structure.
The robust prediction is:
\begin{equation}\label{eq:k0-proportional}
    k_0 \propto m
\end{equation}

The proportionality constant is $\ln 2 \approx 0.693$ under the hypothesis $k_0 = k^*$, but the true constant is empirical.
If $\theta$ is powerful enough to learn modes from very few examples, mode discovery alone drives learning and $k_0 < k^*$.
If $\theta$ requires many repeated observations to characterize a mode, $k_0 > k^*$.

\subsection{Mode Count Estimator}

If the hypothesis holds, the fitted $k_0$ yields a mode count estimator:
\begin{equation}\label{eq:mode-count-estimator}
    \hat{m} = \frac{k_0}{\ln 2} \approx 1.443\, k_0
\end{equation}

This assumes uniform mode weights.
When $k_0 < 1$ (the sigmoid degenerates to exponential decay), the interpretation is that $m$ is small enough that $\theta$ learns all modes within the first few responses, and the curve shape provides no useful mode count estimate.

\subsection{Non-Uniform Mode Weights}

For non-uniform weights $\{w_i\}_{i=1}^m$, the expected probability that draw $k+1$ hits a new mode is:
\begin{equation}
    p_{\mathrm{new}}(k) = \sum_{i=1}^{m} w_i \left(1 - w_i\right)^{k}
\end{equation}

This no longer has a clean closed form for $k^*$.
However, the scaling is governed by the effective mode count.
For weights that are not too extreme, the crossover occurs at roughly:
\begin{equation}\label{eq:kstar-nonuniform}
    k^* \approx \frac{\ln 2}{\sum_i w_i^2}
\end{equation}

where $1/\sum_i w_i^2$ is the inverse Simpson index (a standard measure of effective diversity).
When weights are uniform, $\sum_i w_i^2 = 1/m$ and this recovers $k^* = m \ln 2$.
When weights are concentrated on a few dominant modes, $k^*$ is smaller because the dominant modes are observed early.

\subsection{Estimating Mode Count from the Curve}\label{sec:mode-count-from-curve}

A natural idea is to read the number of effective modes from the \textbf{elbow of the $a_k$ curve}: while $\theta$ is still discovering new modes ($k \lesssim m$), each response is likely novel and $a_k$ remains near $a_1$; once $\theta$ has seen most modes ($k \gtrsim m$), $a_k$ drops toward $a_\infty$.
The transition occurs at $k \approx m$.
The fitted $k_0$ from Eq.~\ref{eq:sigmoid-fit} yields a mode count estimator $\hat{m} = k_0 / \ln 2 \approx 1.443\, k_0$ (assuming uniform mode weights).

However, this prescription requires the curve to be sigmoidal with a visible plateau, which in turn requires modes to be approximately independent.
As shown in Section~\ref{sec:cross-mode-learning}, capable base models exhibit cross-mode learning that violates this assumption: all curves are exponential regardless of $m$, and the fitted $k_0$ hits the optimizer's lower bound ($-10$) at every mode count.
The elbow-based mode count estimator therefore does not work for capable models.
It remains valid for weaker models like GPT-2, but we do not recommend it as a general-purpose tool.

\subsection{Caveats}

\begin{enumerate}
    \item The linear scaling $k_0 \propto m$ is robust (it follows from the combinatorial structure of the coupon collector process).
The constant of proportionality $\ln 2$ is a hypothesis to be tested empirically.

    \item The argument assumes that the sigmoid is a good fit to the $\bar{a}_k$ curve.
When $m$ is small ($m \ll n$), the curve may be better described by exponential decay (the sigmoid with $k_0 < 1$), and the inflection-point interpretation does not apply.

    \item For non-uniform weights, the inverse Simpson index approximation is heuristic.
Modes with very small weight $w_i \lesssim 1/n$ may never appear in the sample and do not contribute to $k^*$ or $k_0$.
\end{enumerate}

